% =====================================================================
%  USER-EDITABLE VARIABLES
%  ---------------------------------------------------------------------
%  Fill in the values below with the data of your thesis.
%  Edit only the text inside the braces { }; do NOT rename the commands,
%  as they are used by the cover page and by the PDF metadata.
% =====================================================================

% --- Academic information --------------------------------------------
% Name of your degree / programme.
\newcommand{\degree}{Bachelor in Computer Science and Engineering}
% Academic year (e.g. 2024-2025).
\newcommand{\academicYear}{2024-2025}
% Type of work, shown on the cover and in the PDF metadata.
\newcommand{\thesisType}{Bachelor's Thesis}

% --- Work information -------------------------------------------------
% Title of the thesis.
\newcommand{\thesisTitle}{MY TITLE}
% Full name of the author.
\newcommand{\thesisAuthor}{AUTHOR}
% Full name of the supervisor / advisor.
\newcommand{\thesisAdvisor}{ADVISOR}
% Place and date of the defence (e.g. Leganes, June 2025).
\newcommand{\placeAndDate}{PLACE, MONTH YEAR}

% --- Institution -----------------------------------------------------
% University name (used in the PDF metadata).
\newcommand{\university}{Universidad Carlos III de Madrid}

% --- Keywords --------------------------------------------------------
% Shown in the Abstract / Resumen and embedded in the PDF metadata.
% Separate the terms with semicolons.
\newcommand{\keywordsEn}{Keyword 1; Keyword 2; Keyword 3}
\newcommand{\keywordsEs}{Palabra clave 1; Palabra clave 2; Palabra clave 3}

% --- Configuration ---------------------------------------------------
% License applied to the work. Allowed values: cc / reserved
\newcommand{\licenseType}{cc}
% Bibliography style. Allowed values: apa / ieee
\newcommand{\bibStyle}{ieee}
% Show the demo chapter (chapters/00_examples) with the template examples. Set it
% to false for the final version (it is just example content): this hides it
% without having to comment anything out in tfg_main.tex.
\newcommand{\showDemoChapter}{true}
% (The PDF/A switch \pdfaOutput lives at the top of tfg_main.tex, because
%  \DocumentMetadata must run before everything else.)
% (The template "Template.*" metadata lives in tfg_template_info.tex; it is
%  maintained by whoever updates the template, you need not touch it for your thesis.)


% =====================================================================
%  GENERATIVE AI (GenAI) USE DECLARATION (mandatory annex)
%  ---------------------------------------------------------------------
%  Rendered by others/ai_declaration.tex. true/false values select and
%  highlight the matching option in the form.
% =====================================================================

% Show the guidance, examples and use-type guide (true) for a working/draft
% version, or hide everything non-essential (hints, comments, unused
% sections) for a production-ready version (false). Switch to false before
% submitting.
\newcommand{\aiShowExplanations}{true}

% Did you use Generative AI in your thesis? (true/false)
% If false, parts 1-3 are omitted and "NO" is marked.
\newcommand{\aiUsed}{true}

% --- Part 1: legal, ethical and responsible behaviour ----------------
% For each question: true = the YES option, false = the NO option.
\newcommand{\aiConfidentialData}{false}    % used confidential data?
\newcommand{\aiCopyrightedMaterial}{false} % used copyrighted material?
\newcommand{\aiPersonalData}{false}        % used personal data?
\newcommand{\aiTermsRespected}{true}       % respected terms of use and ethics?

% --- Part 2: technical use -------------------------------------------
% Name(s) and version(s) of the GenAI system(s)/tool(s) used.
\newcommand{\aiSystemName}{ChatGPT 4o}
% A worked declaration (copy the block in others/ai_declaration.tex for more).
\newcommand{\aiDeclPurpose}{search for information}
\newcommand{\aiDeclPrompt}{Write here the prompt you used}
\newcommand{\aiDeclInteraction}{describe what you did with the AI's answer}

% --- Part 3: reflection on usefulness --------------------------------
\newcommand{\aiReflection}{Write here your personal assessment (free format): the strengths and weaknesses you found when using GenAI tools, and whether they helped your learning, development or conclusions.}


% =====================================================================
%  COLOURS (advanced)
%  ---------------------------------------------------------------------
%  Colour values used throughout the template. Most users will not need
%  to change these. Keep the UC3M corporate blue unless your faculty
%  allows otherwise. Use "R,G,B" (0-255) for RGB and a 6-digit hex for HTML.
%  To ADD a new colour you change TWO files: define its value here, and add a
%  matching \definecolor{<name>}{<model>}{\<yourvalue>} in tfg_uc3m.sty.
% =====================================================================
\newcommand{\colorBrand}{0,0,102}       % UC3M corporate blue (cover, headings) [RGB]
\newcommand{\colorHighlight}{FFFF00}    % text highlight [HTML]
\newcommand{\colorLink}{0000EE}         % hyperlinks, only in PDF/A mode [HTML]
% AI declaration annex
\newcommand{\colorAiSelected}{FFE066}   % selected option [HTML]
\newcommand{\colorAiFill}{D7E9FF}       % values you fill in [HTML]
\newcommand{\colorAiHint}{1A7F5A}       % guidance text [HTML]
% Source-code listings. Greys go 0 = black .. 1 = white; the rest are HTML.
\newcommand{\colorCodeBg}{0.97}         % background (grey)
\newcommand{\colorCodeComment}{0.45}    % comments (grey)
\newcommand{\colorCodeNumbers}{0.55}    % line numbers (grey)
\newcommand{\colorCodeKeyword}{0033CC}  % keywords [HTML]
\newcommand{\colorCodeString}{008000}   % strings [HTML]
